\honest{Serious Stories}{Eliezer Yudkowsky}{2009}

Every Utopia ever built, in philosophy, fiction or religion, has been ``to one degree or another'' a place where you would not want to live. Orwell said so too. \nb{He did: ``All `favourable' Utopias seem to be alike in postulating perfection while being unable to suggest happiness.''} Books on writing say that stories need ``conflict.'' The plainer ones say more: ``Stories are about people's pain'' (Orson Scott Card); ``Every scene must end in disaster'' (Jack Bickham). \nb{Neither is an exact quotation that could be found, but both match the manuals. Card advises that ``your eye should be drawn toward pain''; Bickham's ``disaster'' is a term of art for a setback to the scene's goal, which ``does not often denote an earthquake, a flood, a plane crash.''}

When I was young I thought a positive Singularity, where every problem is wished away, would be ``the end of all stories.'' Later I wondered whether ``no stories'' was a sign that I would not want to live there. So I tried to contemplate the argument: a world where nothing goes wrong has no stories worth reading; a world you would not want to read about is one you would not want to live in; so ``a little pain must fall.''

The great stories, from the \textsc{Iliad} to \textsc{Romeo and Juliet}, are all tragic. \nb{The list is the author's. Aristotle agreed that the unhappy ending is ``the right ending'' for tragedy, and noted that the happy one was preferred ``because of the weakness of the spectators.''} We prefer pleasure in life and prize pain in stories. It is a puzzle. \nb{An old one: Hume's ``Of Tragedy'' (1757) asks how a tragedy pleases when the pleasure seems to depend on sorrow and terror.} Card taught me to make my characters hurt, and Bickham to end every scene in disaster. But ``You simply don't optimize a story the way you optimize a real life,'' and people live well for years without being shot at by assassins. Still, stories where everything goes right grate, as with authors who grow old or successful and stop hurting their characters. Would a life without pain grate the same way? Perhaps stories need only striving, not pain. \nb{Card's own reason for choosing a character in pain is that such a character ``wants things to change. He's likely to act.''}

Or perhaps pain is simply easier to produce. I once read that ``intense sensation'' is a BDSM euphemism for pain, and nothing ordinary makes a person feel as good as the Inquisition's instruments made people hurt. \nb{Baumeister and colleagues (2001) reviewed many studies and concluded that ``bad is stronger than good.''} This goes deeper than loss aversion, under which a loss hurts ``between 2 and 2.5 times'' as much as a gain. \nb{A 2024 meta-analysis puts the mean at 1.955.} Stories must shout; perhaps happiness cannot shout loud enough.

Could we delete pain, and delete adaptation, so that every roast goose tastes as it does after starving? \nb{Adaptation is real but partial. Lottery winners took ``significantly less pleasure from a series of mundane events,'' and later work found that people do not adapt back to neutral.} The brain has ``distinct neural circuits'' for pleasure and pain, and people born without pain do not, so far as I have heard, find too little pleasure intolerable. \nb{A 2008 review stressed how much the circuits overlap. And people born without pain suffer bitten tongues and fractures, and one of the first studied died young, a larger cost than the post's self-inspection.} I would guess pain could be replaced by an urge; I would guess adaptation could be removed; I would guess ``yes, it can be done.'' So David Pearce is ``very probably right'' that suffering can be abolished. \nb{The ``very probably'' rests on those three guesses.}

Is that what we want? Leaving the world as it is will not do. The alternative I like is ``the path of courage'': remove the pain that destroys people, make people stronger, keep broken legs and broken hearts. I admit I may just be ``afraid'' of a very different world. And a child raised without AIDS or slavery might go on to remove heartbreak too. ``I don't know.''
